\documentclass[10pt,a4paper]{amsart}
\usepackage[margin=19mm]{geometry}
\usepackage{amsmath,amssymb,mathtools}
\usepackage[scr=rsfs,cal=euler]{mathalfa}
\usepackage{xfrac}
\usepackage[unicode,hidelinks,bookmarks=false]{hyperref}
\newcommand{\ie}{i.e.\ }
\newcommand{\cf}{cf.\ }
\newcommand{\resp}{respectively\ }
\newcommand{\Subsection}{\subsection}
\providecommand{\nobreakdash}{\nobreak\hbox{-}}
\setlength{\emergencystretch}{2em}
\multlinegap0pt
\usepackage{amssymb}
\usepackage{xcolor}
\newtheorem{prop}{Proposition}
\newcommand{\bN}{\mathbb N}
\title{\bf On image ideals of nice and quasi-nice derivations on $R[X_1,X_2]$}
\author{Nikhilesh Dasgupta$^{*}$ and Animesh Lahiri$^{*}$}
\date{Source: arXiv:2302.13787v6 (CC BY 4.0)}
\begin{document}
\maketitle

\noindent\emph{Source and licence.} \bf On image ideals of nice and quasi-nice derivations on $R[X_1,X_2]$. Version arXiv:2302.13787v6 (CC BY 4.0), distributed by arXiv under the Creative Commons Attribution 4.0 International licence, http://creativecommons.org/licenses/by/4.0/, as recorded in the arXiv licence metadata for this exact version and shown on the abstract page https://arxiv.org/abs/2302.13787v6. Full licence notice: \texttt{licenses/arxiv-2302.13787-CC-BY-4.0.txt}.\par\smallskip

\noindent\textit{Editorial scope.} Counted unit: the article's prop ideal (source0.tex lines 451--460). The article's complete original proof of it is delivered with it (source0.tex lines 461--492, one \texttt{proof} environment of 32 lines). No mathematics is written, repaired or re-derived: the statement and the proof are the article's own lines, in the article's own notation, and no retained line is substituted or re-flowed.  No further source-local context is needed: every cross-reference of every retained line resolves inside the two delivered spans.  Nothing was omitted from any retained span, so the package carries no omission record.  No retained line contains a citation, so the wrapper carries no bibliography and no bibliography entry.  No retained line contains a figure environment, an image-inclusion command or drawing code, so no image is delivered and the package declares no asset dependency.  Editorial formatting only, in addition: an AMS \texttt{amsart} preamble that restates the article's own declarations and macros used by the retained lines, namely the environments \texttt{prop} on separate counters, together with the standard packages those lines need.  The retained environments print in this document as Proposition 1 in order of appearance, whereas the article prints the same items with the numbering of its own full document.  The complete native source of the article as distributed in the arXiv e-print for this exact version is bundled unchanged as \texttt{sources/137912/source0.tex}.
\begin{prop}\label{top degree ideal}
Let $\lambda$ be a weighted degree map on $B~(=R^{[n]})$ and $J=(f_1,\dots,f_m)$ an ideal of $B$. Suppose the following conditions hold:
\begin{enumerate}
	\item [\rm(i)]$\widetilde{f_i}=f_i$ for $1\leqslant i\leqslant m-1$,
	\item[\rm(ii)] $\widetilde{f_m}$ is a non-zerodivisor in $B/(f_1,\dots,f_{m-1})$.
\end{enumerate}

\noindent
Then $\widetilde{J}=(f_1,\dots,f_{m-1}, \widetilde{f_m})$.
\end{prop}

\begin{proof}
Since for each $i$, $1\leqslant i\leqslant m$, $\widetilde{f_i}\in\widetilde{J}$, we clearly have $(f_1,\dots,f_{m-1}, \widetilde{f_m})\subseteq \widetilde{J}$.

\medskip\noindent
Let $I=\{1,\dots,m\}$ and $0\neq g={\underset{i\in I}\sum}r_if_i$, where $r_i\in B$ for all $i\in I$. Let $J_m:=(f_1,\dots,f_{m-1})$ and for each $i\in I$, $\lambda(r_i)=u_i$ and $\lambda(f_i)=v_i$. For $f\in B$, by $f^{(l)}$ we will denote the $l$-degree homogeneous component of $f$ and for $0\leqslant s_i\leqslant u_i$, $0\leqslant t_i\leqslant v_i$, we define $x(i;s_i;t_i):=r_i^{(u_i-s_i)}f_i^{(v_i-t_i)}$. Then $g={\underset{0\leqslant t_i\leqslant v_i}{\underset{0\leqslant s_i\leqslant u_i}{\underset{i\in I}\sum}}}x(i;s_i;t_i).$

\smallskip\noindent Let $M=\{x(i;s_i;t_i):i\in I,0\leqslant s_i\leqslant u_i,0\leqslant t_i\leqslant v_i\}$, $\mu_0:={\rm Max}\{\lambda(x):x\in M \}$ and $M_0:=\{x\in M:\lambda(x)=\mu_0\}$. Clearly, if $x(i;s_i;t_i)\in M_0$, then $s_i=0=t_i$. 

\smallskip\noindent If ${\underset{x\in{M_0}}\sum}x\neq 0$, then $\widetilde {g}={\underset{x\in{M_0}}\sum}x={\underset{M_0}\sum}\widetilde{r_i}\widetilde{f_i}\in (J_m,\widetilde{f_m}).$ If ${\underset{x\in{M_0}}\sum}x=0$, then we will define $\mu_1$ and $M_1$ in the following way:
$$\mu_1:={\rm Max}\{\lambda(x):x\in M\setminus M_0 \} \text{ and }M_1:=\{x\in M\setminus M_0:\lambda(x)=\mu_1\}.$$ %Clearly, if $x(i;s_i;t_i)\in M_1$, then $0\leqslant s_i+t_i\leqslant 1$.
If ${\underset{x\in{M_1}}\sum}x\neq 0$, then $\widetilde {g}={\underset{x\in{M_1}}\sum}x$. In this situation we note that  for some $t_i>0$, $x(i;s_i;t_i)\in M_1$ will imply $i=m$. Clearly, $x(m;s_m;t_m)\in M_1$ for some $t_m>0$ only if $x(m;s_m,0)\in M_0$ and hence $s_m=0$. Since ${\underset{x\in{M_0}}\sum}x= 0$, by hypothesis (ii) of the proposition, $x(m;0;0)\in M_0$ only if $r_m^{(u_m)}\in J_m$. So, for some $t_m>0$, $x(m;s_m;t_m)\in M_1$ only if $x(m;s_m;t_m)~\big{(}=x(m;0;t_m)=r_m^{(u_m)}f_m^{(v_m-t_m)}\big{)}\in J_m$. Hence, ${\underset{x\in{M_1}}\sum}x\neq 0$ will imply $\widetilde{g}\in (J_m,\widetilde{f_m})$.

\medskip\noindent
If ${\underset{M_1}\sum}x=0$, then we will define $\mu_2$ and $M_2$ in the following way and proceed similarly.
$$\mu_2:={\rm Max}\{\lambda(x):x\in (M\setminus M_0)\setminus M_1 \}$$ $$\text{and }M_2:=\{x\in (M\setminus M_0)\setminus M_1:\lambda(x)=\mu_2\}.$$ 
Now, for each positive integer $N$ we consider the following statement:
$$P(N):\text{ If }{\underset{x\in M_j}\sum}x=0 \text{ for all } j<N, \text{ then, for }t_m>0, x(m;s_m;t_m)\in M_N\text{ only if }x(m;s_m;t_m)\in J_m.$$

\noindent
We have already observed that $P(1)$ is true. Assume that $P(1),P(2),\dots,P(N)$ is true for some $N\geqslant 1$. We will show that $P(N+1)$ is also true. Let ${\underset{x\in M_j}\sum}x=0$ for all $j\leqslant N$ and $x(m;s_m^{'};t_m^{'})\in M_{N+1}$ for some fixed $s_m^{'},t_m^{'}$ with $0\leqslant s_m^{'}\leqslant u_m$ and $1\leqslant t_m^{'}\leqslant v_m$. Then $x(m;s_m^{'};0)\in M_{j_0}$ for some $0\leqslant j_0\leqslant N$. We observe the following:
\begin{enumerate}
	\item $0={\underset{M_{j_0}}\sum}x(i;s_i;t_i)={\underset{i\neq m}{\underset{M_{j_0}}\sum}}x(i;s_i;0)+{\underset{M_{j_0}}\sum}x(m;s_m;0)+{\underset{t_m>0}{\underset{M_{j_0}}\sum}}x(m;s_m;t_m)$,
	\item ${\underset{i\neq m}{\underset{M_{j_0}}\sum}}x(i;s_i;0)\in J_m$,
	\item Since $P(j_0)$ is true, ${\underset{t_m> 0}{\underset{M_{j_0}}\sum}}x(m;s_m;t_m)\in J_m$,
	\item Since $x(m;s_m^{'};0)\in M_{j_0}$, $x(m;s_m;0)\in M_{j_0}$ if and only if $s_m=s_m^{'}$.
\end{enumerate}
\medskip\noindent 
Hence $x(m,s_m^{'},0)\in J_m$ and  by hypothesis (ii), $r_m^{(u_m-s_m^{'})}\in J_m$. So, $x(m;s_m^{'};t_m^{'})\in J_m$. Thus, $P(n+1)$ is true and hence $P(n)$ is true for all $n\in\bN$.

\medskip\noindent
Now, for $n\in\bN$, if ${\underset{x\in M_j}\sum}x= 0$ for all $j<n$ and ${\underset{x\in M_n}\sum}x\neq 0$, then $\widetilde{g}={\underset{x\in M_n}\sum}x\in(J_m,\widetilde{f_m})$. If ${\underset{x\in M_j}\sum}x=0$ for all $j\leqslant n$, then we will consider $\mu_{n+1},M_{n+1}$ and repeat the arguments for $M_{n+1}$. Since $M$ is a finite set this process will stop after a finite number of steps.
\end{proof}

\end{document}
